\documentclass[10pt,a4paper]{amsart}
\usepackage[margin=19mm]{geometry}
\usepackage{amsmath,amssymb,mathtools}
\usepackage[scr=rsfs,cal=euler]{mathalfa}
\usepackage{xfrac}
\usepackage[unicode,hidelinks,bookmarks=false]{hyperref}
\newcommand{\ie}{i.e.\ }
\newcommand{\cf}{cf.\ }
\newcommand{\resp}{respectively\ }
\newcommand{\Subsection}{\subsection}
\providecommand{\nobreakdash}{\nobreak\hbox{-}}
\setlength{\emergencystretch}{2em}
\multlinegap0pt
\usepackage{bm}
\usepackage{algorithm}\usepackage{algpseudocode}
\usepackage{tikz}
\usepackage{amssymb}
\usepackage[shortlabels]{enumitem}
\newtheorem{proposition}{Proposition}
\title{Optimal Selection with Balanced Market Share: Static and Dynamic Assortment Optimization}
\author{Omar El Housni, Qing Feng, Huseyin Topaloglu}
\date{Source: arXiv:2507.05606v2 (CC BY 4.0)}
\begin{document}
\maketitle

\noindent\emph{Source and licence.} Optimal Selection with Balanced Market Share: Static and Dynamic Assortment Optimization. Version arXiv:2507.05606v2 (CC BY 4.0), distributed by arXiv under the Creative Commons Attribution 4.0 International licence, http://creativecommons.org/licenses/by/4.0/, as recorded in the arXiv licence metadata for this exact version and shown on the abstract page https://arxiv.org/abs/2507.05606v2. Full licence notice: \texttt{licenses/arxiv-2507.05606-CC-BY-4.0.txt}.\par\smallskip

\noindent\textit{Editorial scope.} Counted unit: the article's proposition prop:heuristic (source0.tex lines 2182--2184). The article's complete original proof of it is delivered with it (source0.tex lines 2186--2235, one \texttt{proof} environment of 50 lines). No mathematics is written, repaired or re-derived: the statement and the proof are the article's own lines, in the article's own notation, and no retained line is substituted or re-flowed.  Retained uncounted as the source-local context the proof needs: lines 666--673; lines 681--690; lines 1942--1951.  Nothing was omitted from any retained span, so the package carries no omission record.  No retained line contains a citation, so the wrapper carries no bibliography and no bibliography entry.  No retained line contains a figure environment, an image-inclusion command or drawing code, so no image is delivered and the package declares no asset dependency.  Editorial formatting only, in addition: an AMS \texttt{amsart} preamble that restates the article's own declarations and macros used by the retained lines, namely the environments \texttt{proposition} on separate counters, together with the standard packages those lines need.  The retained environments print in this document as Proposition 1 in order of appearance, whereas the article prints the same items with the numbering of its own full document.  The complete native source of the article as distributed in the arXiv e-print for this exact version is bundled unchanged as \texttt{sources/181190/source0.tex}.
\begin{equation}\label{prob:inventory}
\tag{\texttt{BMS-Dynamic}}
\begin{aligned}
&\max_{\pi\in\Pi}&&{\sum_{i\in\mathcal{N}}r_i\mathbb{E}[X_{iT}^\pi]}\\
&\text{s.t.}&&{\mathbb{P}(X_{iT}^\pi\leq c_i)=1,\ \forall i\in\mathcal{N},}\\
&&&{\mathbb{E}[X_{iT}^\pi]\in\{0\}\cup \Big[\alpha\cdot \max_{j\in\mathcal{N}}\mathbb{E}[X_{jT}^\pi],\infty\Big),\ \forall i\in\mathcal{N}.}
\end{aligned}
\end{equation}

\begin{equation}\label{prob:inventory_upper_bound}
\tag{\texttt{Upper-Bound}}
\begin{aligned}
&\max_{\bm{x}}&&{T\sum_{i\in\mathcal{N}}r_ix_i}\\
&\text{s.t.}&&{x_0+\sum_{i\in\mathcal{N}}x_i=1},\\
&&&{0\leq x_i\leq v_ix_0,\ \forall i\in\mathcal{N}},\\
&&&{x_i\leq c_i/T,\ \forall i\in\mathcal{N}},\\
&&&{x_i\in\{0\}\cup\Big[\alpha\cdot\max_{j\in\mathcal{N}}x_j,\infty\Big),\ \forall i\in\mathcal{N}.}
\end{aligned}
\end{equation}

\begin{equation}\label{prob:resolving}
\tag{\texttt{BMS-H}}
\begin{aligned}
&\max_{\bm{x}}&&{\sum_{i\in\bar{S}}r_ix_i}\\
&\text{s.t.}&&{x_0+\sum_{i\in\bar{S}}x_i=1},\\
&&&{0\leq x_i\leq v_ix_0,\ \forall i\in\bar{S}},\\
&&&\sum_{s=1}^{t-1}\tilde{x}_i^{(s)}+(T-t+1)x_i\geq \alpha\cdot\max_{j\in\mathcal{N}}\sum_{s=1}^{t-1}\tilde{x}_j^{(s)}+(T-t+1)x_j,\ \forall i\in\bar{S},\\
&&&{(T-t+1)x_{i}\leq c_i-X_{i,t-1},\ \forall i\in\bar{S}.}
\end{aligned}
\end{equation}

\begin{proposition}\label{prop:heuristic}
For both \textsc{Heuristic 1} and \textsc{Heuristic 2}, \ref{prob:resolving}~is always feasible when resolved, and both heuristics are feasible to~\ref{prob:inventory}.
\end{proposition}

\begin{proof}
We denote the set of all possible history of sales up till time period $t$ as $\mathcal{H}_t$. Recall that $\tilde{\bm{x}}^{(t)}$ gives the purchase probabilities at time period $t$. Since $\tilde{\bm{x}}^{(t)}$ is determined by the history of sales up till time period $t-1$, we consider $\tilde{\bm{x}}^{(t)}$ as a random variable in the history of sales space $\mathcal{H}_{t-1}$. As a generalization of both \textsc{Heuristic 1} and \textsc{Heuristic 2}, we consider any policy $\pi$ in the following form: the policy always resolves~\ref{prob:resolving}~whenever some product runs out of inventory exactly in the previous time period, but the policy may also resolve~\ref{prob:resolving}~at other time periods depending on the history of sales.

Recall that $\bar{S}=\{i\in\mathcal{N}:x_i^*>0\}$, where $\bm{x}^*$ is an optimal solution to the \ref{prob:inventory_upper_bound} problem. We first prove that in any of the aforementioned policies, for any $t\in\{1,2,\dots,T\}$, \ref{prob:resolving} is always feasible up till time period $t$, and for all $i,j\in\bar{S}$, the following inequality holds with probability one,
\begin{equation}\label{eq:balance_t}
\sum_{s=1}^t \tilde{x}_i^{(s)}\geq \alpha\sum_{s=1}^t\tilde{x}_j^{(s)}.
\end{equation}
We prove the result by induction. Letting $\bm{x}^*$ be an optimal solution to the~\ref{prob:inventory_upper_bound}~problem, since by definition $\tilde{\bm{x}}^{(1)}=\bm{x}^*$, it is easy to verify that when $t=1$, \ref{prob:resolving}~is always feasible, and by the feasibility of $\bm{x}^*$ we have $\tilde{\bm{x}}^{(1)}$ satisfies $\tilde{x}_i^{(1)}\geq \alpha\tilde{x}_j^{(1)}$ for all $i,j\in\bar{S}$, therefore the result holds for $t=1$. 

Suppose \ref{prob:resolving} is always feasible up till time step $t_0$, and \eqref{eq:balance_t} holds for all $t\leq t_0$ for some $t_0<T$. Consider any history of sales $\bm{H}_{t_0}\in\mathcal{H}_{t_0}$ up till time period $t_0$. Suppose under history of sales $\bm{H}_{t_0}$, the policy resolves~\ref{prob:resolving} at time period $t_0+1$. By the induction assumption we have that for all $i,j\in\bar{S}$, $\sum_{s=1}^{t_0}\tilde{x}_i^{(s)}\geq \alpha\sum_{s=1}^{t_0}\tilde{x}_j^{(s)}$. Therefore when resolving~\ref{prob:resolving}~at time period $t_0+1$, it is easy to verify that $x_i=0$ for all $i\in\bar{S}$, is feasible to~\ref{prob:resolving}, therefore~\ref{prob:resolving}~is always feasible up till time period $t_0+1$.

Then we prove that for all $i,j\in\bar{S}$, we have $\sum_{s=1}^{t_0+1}\tilde{x}_i^{(s)}\geq \alpha\sum_{s=1}^{t_0+1}\tilde{x}_j^{(s)}$. Consider any history of sales $\bm{H}_{t_0}\in\mathcal{H}_{t_0}$. We define $\tau$ as the last time period up till $t_0+1$ in which the policy resolves~\ref{prob:resolving} under history of sales $\bm{H}_{t_0}$. If the policy resolves~\ref{prob:resolving}~at time period $t_0+1$ then we define $\tau=t_0+1$. By the induction assumption we have that for all $i,j\in\bar{S}$
\begin{equation}\label{eq:resolving_1}
\sum_{s=1}^{\tau-1}\tilde{x}_i^{(s)}\geq \alpha\sum_{s=1}^{\tau-1}\tilde{x}_j^{(s)}.
\end{equation}
Since $\tilde{\bm{x}}^{(\tau)}$ is an optimal solution when resolving~\ref{prob:resolving}~at time period $\tau$, by the third set of constraints in~\ref{prob:resolving}~we get
\begin{equation}\label{eq:resolving_2}
\sum_{s=1}^{\tau-1}\tilde{x}_i^{(s)}+(T-\tau+1)\tilde{x}_i^{(\tau)}\geq \alpha\sum_{s=1}^{\tau-1}\tilde{x}_j^{(s)}+\alpha(T-\tau+1)\tilde{x}_j^{(\tau)}.
\end{equation}
By definition of $\tau$ we have that no resolving happens between time period $\tau+1$ and $t_0+1$, thus $\tilde{\bm{x}}_i^{(s)}=\tilde{\bm{x}}^{(\tau)}$ for all $\tau\leq s\leq t_0+1$. Therefore $\sum_{s=1}^{t_0+1}\tilde{x}_i^{(s)}=\sum_{s=1}^{\tau-1}\tilde{x}_i^{(s)}+(t_0-\tau+2)\tilde{x}_i^{\tau}$ for all $i\in\bar{S}$. Combining \eqref{eq:resolving_1} and \eqref{eq:resolving_2} we have that for all $i,j\in\bar{S}$,
\begin{align*}
&\sum_{s=1}^{t_0+1}\tilde{x}_i^{(s)}=\sum_{s=1}^{\tau-1}\tilde{x}_i^{(s)}+(t_0-\tau+2)\tilde{x}_i^{(\tau)}=\dfrac{T-t_0-1}{T-\tau+1}\sum_{s=1}^{\tau-1}\tilde{x}_i^{(s)}+\dfrac{t_0-\tau+2}{T-\tau+1}\Big(\sum_{s=1}^{\tau-1}\tilde{x}_i^{(s)}+(T-\tau+1)\tilde{x}_i^{(\tau)}\Big)\\
\geq&\dfrac{T-t_0-1}{T-\tau+1}\alpha\sum_{s=1}^{\tau-1}\tilde{x}_j^{(s)}+\dfrac{t_0-\tau+2}{T-\tau+1}\alpha\Big(\sum_{s=1}^{\tau-1}\tilde{x}_j^{(s)}+(T-\tau+1)\tilde{x}_j^{(\tau)}\Big)=\alpha\sum_{s=1}^{t_0+1}\tilde{x}_j^{(s)}.
\end{align*}
Here the inequality is due to $t_0<T$ and $\tau\leq t_0+1$, thus $T-t_0-1\geq 0$ and $t_0-\tau+2\geq 0$. Therefore we conclude that the result holds for $t=t_0+1$. By induction \ref{prob:resolving} is always feasible, and \eqref{eq:balance_t} holds for all $t\in\{1,2,\dots,T\}$. Therefore~\ref{prob:resolving}~is feasible whenever resolved. We further have $\sum_{t=1}^T\tilde{x}_i^{(t)}\geq \alpha\sum_{t=1}^T\tilde{x}_j^{(t)}$ holds with probability one for all $i,j\in\bar{S}$.

We claim that
\begin{equation*}
\mathbb{E}[X_{iT}^\pi]=\sum_{t=1}^T\mathbb{E}[\tilde{x}_i^{(t)}].
\end{equation*}
Consider any $t\in\{1,2,\dots,T\}$ and any $i\in\bar{S}$, we have that
\begin{equation*}
\mathbb{E}[X_{it}^\pi-X_{i,t-1}^\pi|\bm{H}_{t-1}]=\tilde{x}_i^{(t)}.
\end{equation*}
Taking the expectation over $\bm{H}_{t-1}$ we get
\begin{equation*}
\mathbb{E}[X_{it}^\pi-X_{i,t-1}^\pi]=\mathbb{E}[\tilde{x}_i^{(t)}].
\end{equation*}
Since $X_{i0}=0$ for all $i\in\bar{S}$, we get
\begin{equation*}
\mathbb{E}[X_{iT}^\pi]=\sum_{t=1}^T\mathbb{E}[X^\pi_{it}-X^\pi_{i,t-1}]=\sum_{t=1}^T\mathbb{E}[\tilde{x}_i^{(t)}].
\end{equation*}
We have proven that $\sum_{t=1}^T\tilde{x}_i^{(t)}\geq \alpha\sum_{t=1}^T\tilde{x}_j^{(t)}$ holds with probability one for all $i,j\in\bar{S}$. Thus for all $i,j\in\bar{S}$,
$$\mathbb{E}[X_{iT}^\pi]=\sum_{t=1}^T\mathbb{E}[\tilde{x}_i^{(t)}]\geq\alpha\sum_{t=1}^T\mathbb{E}[\tilde{x}_j^{(t)}]=\alpha\mathbb{E}[X_{jT}^\pi].$$
It is easy to verify that $\mathbb{E}[X_{iT}^\pi]=0$ if $i\notin\bar{S}$, therefore for any $i\in\mathcal{N}$, we get
\begin{equation*}
\mathbb{E}[X_{iT}^\pi]\in\{0\}\cup\Big[\alpha\cdot\max_{j\in\mathcal{N}}\mathbb{E}[X_{jT}^\pi],\infty\Big).
\end{equation*}
Furthermore, as soon as any product $i\in\bar{S}$ runs out of inventory, we resolve~\ref{prob:resolving}, and the purchase probability for product $i$ is always zero afterwards. Therefore $X_{it}^\pi\leq c_i$ with probability one. Then we conclude that the policy is feasible to~\ref{prob:inventory}. In particular, both \textsc{Heuristic 1} and \textsc{Heuristic 2} are well-defined policies and feasible to~\ref{prob:inventory}.
\end{proof}

\end{document}
